% GENERATED FILE. Edit canonical lessons, then run npm run generate:books.
% canonical-source-hash: fnv1a64:1dee7341c1614b40
% canonical-lessons: ML-C454-unnayikkuka, ML-C454-pratirodhikkuka, ML-C454-sarivekkuka, ML-C454-bodhyappetuka, ML-C454-sthirikarikkuka

\chapter{Raising Questions}
\label{ch:ml-raising-questions}
\hlchaptermodality{454}

\begin{chapteropening}
\textbf{By the end of this chapter:} \emph{I can say to bring up, to defend, to uphold, to be convinced and to confirm.}

Complete the last lesson of chapter 454: I can say to bring up, to defend, to uphold, to be convinced and to confirm.
\end{chapteropening}

\section[\texorpdfstring{\ml{ഉന്നയിക്കുക} (unnayikkuka)}{ഉന്നയിക്കുക (unnayikkuka)}]{\ml{ഉന്നയിക്കുക} (unnayikkuka) — to bring up (an issue)}

\label{lesson:ML-C454-unnayikkuka}



\begin{quote}
Before the new one: say the Malayalam for to question, then the Malayalam for to respond.
\end{quote}

\subsection*{You'll want to know: \ml{ഉന്നയിക്കുക}}
\textbf{\ml{ഉന്നയിക്കുക}} — \emph{unnayikkuka} — \textquotedblleft{}to bring up (an issue)\textquotedblright{}.

\begin{quote}
\ml{യോഗത്തിൽ} \ml{അവൾ} \ml{വെള്ളത്തിന്റെ} \ml{പ്രശ്നം} \ml{ഉന്നയിച്ചു}.
\end{quote}

\begin{quote}
\emph{yōgattil avaḷ veḷḷattinṟe praśnaṁ unnayiccu.}
\end{quote}

\begin{quote}
At the meeting she brought up the water problem.
\end{quote}

Past: \ml{ഉന്നയിച്ചു}, \emph{unnayiccu}.

Said of a question, a demand or a complaint.

\begin{cousinweb}
\textbf{\ml{ഉന്നയിക്കുക}} (\emph{unnayikkuka}) is built from Sanskrit \emph{unnaya}, \textquotedblleft{}the act of leading up, raising\textquotedblright{} (Monier-Williams), with the Malayalam verb ending -\ml{ിക്കുക}.
\end{cousinweb}

\begin{grammarlens}[title={What you've built}]
One more word for saying what happens if.
\end{grammarlens}

\subsection*{Your turn}
\begin{itemize}
\raggedright
  \item \emph{Say it:} \emph{unnayikkuka}
  \item \emph{Say it:} \emph{unnayikkuka}, once more
  \item \emph{Say it:} \emph{yōgattil avaḷ veḷḷattinṟe praśnaṁ unnayiccu.} — At the meeting she brought up the water problem.
  \item \emph{Recall:} say the Malayalam for to question, then the Malayalam for to respond, then say \emph{unnayikkuka} again
\end{itemize}

\begin{tcolorbox}[breakable,colback=teal!4,colframe=teal!35!black,title={Before you move on}]
What does \ml{ഉന്നയിക്കുക} mean? (\textquotedblleft{}To bring up (an issue)\textquotedblright{}.) Say it once more.
\end{tcolorbox}
\section[\texorpdfstring{\ml{പ്രതിരോധിക്കുക} (pratirōdhikkuka)}{പ്രതിരോധിക്കുക (pratirōdhikkuka)}]{\ml{പ്രതിരോധിക്കുക} (pratirōdhikkuka) — to defend; to resist}

\label{lesson:ML-C454-pratirodhikkuka}



\begin{quote}
Before the new one: say the Malayalam for to respond, then the Malayalam for to bring up.
\end{quote}

\subsection*{You'll want to know: \ml{പ്രതിരോധിക്കുക}}
\textbf{\ml{പ്രതിരോധിക്കുക}} — \emph{pratirōdhikkuka} — \textquotedblleft{}to defend; to resist\textquotedblright{}.

\begin{quote}
\ml{അവൻ} \ml{തന്റെ} \ml{തീരുമാനം} \ml{ശക്തമായി} \ml{പ്രതിരോധിച്ചു}.
\end{quote}

\begin{quote}
\emph{avan tanṟe tīrumānaṁ śaktamāyi pratirōdhiccu.}
\end{quote}

\begin{quote}
He firmly defended his decision.
\end{quote}

Past: \ml{പ്രതിരോധിച്ചു}, \emph{pratirōdhiccu}.

The noun is \ml{പ്രതിരോധം}, defence, resistance.

\begin{cousinweb}
\textbf{\ml{പ്രതിരോധിക്കുക}} (\emph{pratirōdhikkuka}) is built from Sanskrit \emph{pratirodha}, \textquotedblleft{}opposition, obstruction\textquotedblright{} (Monier-Williams), with the Malayalam verb ending -\ml{ിക്കുക}.
\end{cousinweb}

\begin{grammarlens}[title={What you've built}]
One more word for saying what happens if.
\end{grammarlens}

\subsection*{Your turn}
\begin{itemize}
\raggedright
  \item \emph{Say it:} \emph{pratirōdhikkuka}
  \item \emph{Say it:} \emph{pratirōdhikkuka}, once more
  \item \emph{Say it:} \emph{avan tanṟe tīrumānaṁ śaktamāyi pratirōdhiccu.} — He firmly defended his decision.
  \item \emph{Recall:} say the Malayalam for to respond, then the Malayalam for to bring up, then say \emph{pratirōdhikkuka} again
\end{itemize}

\begin{tcolorbox}[breakable,colback=teal!4,colframe=teal!35!black,title={Before you move on}]
What does \ml{പ്രതിരോധിക്കുക} mean? (\textquotedblleft{}To defend; to resist\textquotedblright{}.) Say it once more.
\end{tcolorbox}
\section[\texorpdfstring{\ml{ശരിവെക്കുക} (śarivekkuka)}{ശരിവെക്കുക (śarivekkuka)}]{\ml{ശരിവെക്കുക} (śarivekkuka) — to uphold; to endorse}

\label{lesson:ML-C454-sarivekkuka}



\begin{quote}
Before the new one: say the Malayalam for to bring up, then the Malayalam for to defend.
\end{quote}

\subsection*{You'll want to know: \ml{ശരിവെക്കുക}}
\textbf{\ml{ശരിവെക്കുക}} — \emph{śarivekkuka} — \textquotedblleft{}to uphold; to endorse\textquotedblright{}.

\begin{quote}
\ml{മേൽക്കോടതി} \ml{ആ} \ml{വിധി} \ml{ശരിവെച്ചു}.
\end{quote}

\begin{quote}
\emph{mēlkkōṭati ā vidhi śariveccu.}
\end{quote}

\begin{quote}
The higher court upheld that verdict.
\end{quote}

Past: \ml{ശരിവെച്ചു}, \emph{śariveccu}.

\ml{ശരി}, right, with \ml{വെക്കുക}, to put: to set a thing down as right.

\begin{cousinweb}
\textbf{\ml{ശരിവെക്കുക}} (\emph{śarivekkuka}) is a Malayalam compound: \ml{ശരി}, \textquotedblleft{}right, correct\textquotedblright{}, whose origin is uncertain (often traced to Sanskrit \emph{sadṛśa}, \textquotedblleft{}alike\textquotedblright{}, CDIAL 13120), and \ml{വെക്കുക}, \textquotedblleft{}to put\textquotedblright{}, native Dravidian (DEDR 5549).
\end{cousinweb}

\begin{grammarlens}[title={What you've built}]
One more word for saying what happens if.
\end{grammarlens}

\subsection*{Your turn}
\begin{itemize}
\raggedright
  \item \emph{Say it:} \emph{śarivekkuka}
  \item \emph{Say it:} \emph{śarivekkuka}, once more
  \item \emph{Say it:} \emph{mēlkkōṭati ā vidhi śariveccu.} — The higher court upheld that verdict.
  \item \emph{Recall:} say the Malayalam for to bring up, then the Malayalam for to defend, then say \emph{śarivekkuka} again
\end{itemize}

\begin{tcolorbox}[breakable,colback=teal!4,colframe=teal!35!black,title={Before you move on}]
What does \ml{ശരിവെക്കുക} mean? (\textquotedblleft{}To uphold; to endorse\textquotedblright{}.) Say it once more.
\end{tcolorbox}
\section[\texorpdfstring{\ml{ബോധ്യപ്പെടുക} (bōdhyappeṭuka)}{ബോധ്യപ്പെടുക (bōdhyappeṭuka)}]{\ml{ബോധ്യപ്പെടുക} (bōdhyappeṭuka) — to be convinced}

\label{lesson:ML-C454-bodhyappetuka}



\begin{quote}
Before the new one: say the Malayalam for to defend, then the Malayalam for to uphold.
\end{quote}

\subsection*{You'll want to know: \ml{ബോധ്യപ്പെടുക}}
\textbf{\ml{ബോധ്യപ്പെടുക}} — \emph{bōdhyappeṭuka} — \textquotedblleft{}to be convinced\textquotedblright{}.

\begin{quote}
\ml{അവന്റെ} \ml{വിശദീകരണം} \ml{കേട്ടപ്പോൾ} \ml{എനിക്ക്} \ml{ബോധ്യപ്പെട്ടു}.
\end{quote}

\begin{quote}
\emph{avanṟe viśadīkaraṇaṁ kēṭṭappōḷ enikkŭ bōdhyappeṭṭu.}
\end{quote}

\begin{quote}
When I heard his explanation, I was convinced.
\end{quote}

Past: \ml{ബോധ്യപ്പെട്ടു}, \emph{bōdhyappeṭṭu}.

The person convinced takes -\ml{ക്ക്}: \ml{എനിക്ക്} \ml{ബോധ്യപ്പെട്ടു}. Compare \ml{ബോധ്യപ്പെടുത്തുക}, to convince.

\begin{cousinweb}
\textbf{\ml{ബോധ്യപ്പെടുക}} (\emph{bōdhyappeṭuka}) is built from Sanskrit \emph{bodhya}, \textquotedblleft{}to be known or understood\textquotedblright{} (Monier-Williams), with -\ml{പ്പെടുക} from native \ml{പെടുക} (DEDR 3853): to become understood.
\end{cousinweb}

\begin{grammarlens}[title={What you've built}]
One more word for saying what happens if.
\end{grammarlens}

\subsection*{Your turn}
\begin{itemize}
\raggedright
  \item \emph{Say it:} \emph{bōdhyappeṭuka}
  \item \emph{Say it:} \emph{bōdhyappeṭuka}, once more
  \item \emph{Say it:} \emph{avanṟe viśadīkaraṇaṁ kēṭṭappōḷ enikkŭ bōdhyappeṭṭu.} — When I heard his explanation, I was convinced.
  \item \emph{Recall:} say the Malayalam for to defend, then the Malayalam for to uphold, then say \emph{bōdhyappeṭuka} again
\end{itemize}

\begin{tcolorbox}[breakable,colback=teal!4,colframe=teal!35!black,title={Before you move on}]
What does \ml{ബോധ്യപ്പെടുക} mean? (\textquotedblleft{}To be convinced\textquotedblright{}.) Say it once more.
\end{tcolorbox}
\section[\texorpdfstring{\ml{സ്ഥിരീകരിക്കുക} (sthirīkarikkuka)}{സ്ഥിരീകരിക്കുക (sthirīkarikkuka)}]{\ml{സ്ഥിരീകരിക്കുക} (sthirīkarikkuka) — to confirm}

\label{lesson:ML-C454-sthirikarikkuka}



\begin{quote}
Before the new one: say the Malayalam for to uphold, then the Malayalam for to be convinced.
\end{quote}

\subsection*{You'll want to know: \ml{സ്ഥിരീകരിക്കുക}}
\textbf{\ml{സ്ഥിരീകരിക്കുക}} — \emph{sthirīkarikkuka} — \textquotedblleft{}to confirm\textquotedblright{}.

\begin{quote}
\ml{യോഗം} \ml{നാളെയാണെന്ന്} \ml{പ്രിൻസിപ്പൽ} \ml{സ്ഥിരീകരിച്ചു}.
\end{quote}

\begin{quote}
\emph{yōgaṁ nāḷeyāṇennŭ prinsippal sthirīkariccu.}
\end{quote}

\begin{quote}
The principal confirmed that the meeting is tomorrow.
\end{quote}

Past: \ml{സ്ഥിരീകരിച്ചു}, \emph{sthirīkariccu}.

The noun is \ml{സ്ഥിരീകരണം}, confirmation.

\begin{cousinweb}
\textbf{\ml{സ്ഥിരീകരിക്കുക}} (\emph{sthirīkarikkuka}) is built from Sanskrit \emph{sthira}, \textquotedblleft{}firm, hard, durable\textquotedblright{} (CDIAL 13771), with -\ml{ീകരിക്കുക}, \textquotedblleft{}to make into\textquotedblright{}: to make a thing firm.
\end{cousinweb}

\begin{grammarlens}[title={What you've built}]
One more word for saying what happens if.
\end{grammarlens}

\subsection*{Your turn}
\begin{itemize}
\raggedright
  \item \emph{Say it:} \emph{sthirīkarikkuka}
  \item \emph{Say it:} \emph{sthirīkarikkuka}, once more
  \item \emph{Say it:} \emph{yōgaṁ nāḷeyāṇennŭ prinsippal sthirīkariccu.} — The principal confirmed that the meeting is tomorrow.
  \item \emph{Recall:} say the Malayalam for to uphold, then the Malayalam for to be convinced, then say \emph{sthirīkarikkuka} again
\end{itemize}

\begin{tcolorbox}[breakable,colback=teal!4,colframe=teal!35!black,title={Before you move on}]
What does \ml{സ്ഥിരീകരിക്കുക} mean? (\textquotedblleft{}To confirm\textquotedblright{}.) Say it once more.
\end{tcolorbox}
